\section{Results}
\label{sec:results}

We report assessment results under two conditions: cold start (modules freshly initialized with no accumulated state) and warm start (modules exercised through five consciousness cycles with synthetic stimuli before assessment). The cold-start condition serves as a controlled baseline; the warm-start condition demonstrates the effect of even minimal operational history on assessment scores. Cold-start results are perfectly deterministic (repeated runs return identical scores). Warm-start results are reported as means $\pm$ standard deviations across ten independent runs.

All results were obtained on consumer hardware (Apple M4 Pro, 64 GB unified memory). Assessment processing time was under 1 ms in all conditions.

\subsection{Cold-Start Assessment}
\label{sec:cold-start}

With all modules freshly initialized and no prior processing history, the framework scores:

\begin{table}[t]
\centering
\begin{tabular}{ll}
\toprule
Metric & Value \\
\midrule
Overall score & 0.430 \\
Passing indicators & 11/20 (55\%) \\
Architecture functional & No \\
Processing time & $<$ 1 ms \\
\bottomrule
\end{tabular}
\caption{Cold-start assessment summary.}
\label{tab:cold-start-summary}
\end{table}

\textbf{Theory scores (cold start):}

\begin{table}[t]
\centering
\small
\begin{tabular}{llp{8cm}}
\toprule
Theory & Score & Interpretation \\
\midrule
FEP & 0.850 & Strong --- hierarchy exists structurally \\
GWT & 0.450 & Moderate --- workspace has not yet processed candidates \\
HOT & 0.450 & Moderate --- no accumulated higher-order thought history \\
AST & 0.358 & Low --- no tracking history or prediction accuracy \\
BLT & 0.300 & Low --- binding and epistemic depth require inference history \\
IIT & 0.250 & Low --- Φ approximation limited without connectivity data \\
RPT & 0.100 & Low --- requires active substrates with feedback dynamics \\
\bottomrule
\end{tabular}
\caption{Theory scores at cold start.}
\label{tab:cold-start-theory}
\end{table}

Cold-start results are \textbf{perfectly reproducible}: repeated runs return identical scores (0.430, 11/20 passing). This is expected --- with no accumulated state and no random input, the scoring functions are deterministic.

The architecture does not meet the ``functional'' criterion at cold start because the overall score (0.430) falls below the 0.5 threshold. This is by design: a system that scores ``architecture functional'' without any operational history would be measuring structural properties rather than dynamic function.

\subsection{Warm-Start Assessment}
\label{sec:warm-start}

After five consciousness cycles with synthetic stimuli (random workspace candidates with varying salience, emotional charge, and novelty), the framework scores improve substantially. Results across ten independent runs:

\begin{table}[t]
\centering
\begin{tabular}{ll}
\toprule
Metric & Value \\
\midrule
Overall score & $0.506 \pm 0.007$ \\
Passing indicators & 14/20 (70\%) \\
Architecture functional & Yes \\
Processing time & $<$ 1 ms \\
\bottomrule
\end{tabular}
\caption{Warm-start assessment summary (mean $\pm$ std across 10 runs).}
\label{tab:warm-start-summary}
\end{table}

\textbf{Theory scores (warm start, mean $\pm$ std):}

\begin{table}[t]
\centering
\begin{tabular}{llll}
\toprule
Theory & Cold & Warm & Change \\
\midrule
FEP & 0.850 & $0.855 \pm 0.024$ & $+0.005$ \\
GWT & 0.450 & $0.746 \pm 0.020$ & \textbf{+0.296} \\
HOT & 0.450 & $0.450 \pm 0.000$ & --- \\
AST & 0.358 & $0.417 \pm 0.000$ & $+0.059$ \\
BLT & 0.300 & $0.300 \pm 0.000$ & --- \\
IIT & 0.250 & $0.250 \pm 0.000$ & --- \\
RPT & 0.100 & $0.100 \pm 0.000$ & --- \\
\bottomrule
\end{tabular}
\caption{Theory scores: cold start vs.\ warm start.}
\label{tab:cold-warm-comparison}
\end{table}

The most striking change is GWT, which jumps from 0.450 to 0.746 --- a 66\% increase. This is because GWT's indicators (global broadcast, ignition dynamics, global ignition nuanced) measure dynamic properties that only manifest when the workspace has actually processed candidates through competition, ignition, and broadcast. Five cycles of synthetic input are sufficient to demonstrate these dynamics.

\textbf{Indicators that change from FAIL to PASS after warmup:}

\begin{table}[t]
\centering
\begin{tabular}{llll}
\toprule
Indicator & Cold & Warm (mean) & Theory \\
\midrule
global\_broadcast & 0.450 & 1.000 & GWT \\
global\_ignition\_nuanced & 0.350 & $0.725 \pm 0.031$ & GWT \\
ignition\_dynamics & 0.400 & $0.660 \pm 0.049$ & GWT \\
\bottomrule
\end{tabular}
\caption{Indicators that flip from failing to passing after warmup.}
\label{tab:flip-indicators}
\end{table}

Three indicators flip from failing to passing, all belonging to GWT. This is consistent with GWT being the theory most dependent on operational history --- workspace competition is a process, not a structure, and it requires actual candidates to demonstrate.

Warm-start variance is small (overall $\pm 0.007$) and concentrated in the indicators that depend on random synthetic input: ignition\_dynamics, prediction\_error\_minimization, sparse\_smooth\_coding, and global\_ignition\_nuanced. Theory-level scores for HOT, AST, BLT, IIT, and RPT show zero variance because their indicators measure structural properties unaffected by the specific content of warmup stimuli.

\subsection{Per-Indicator Results (Warm Start)}
\label{sec:per-indicator}

The full indicator breakdown reveals which architectural features are present after warmup (mean scores across 5 runs):

\textbf{Passing (14/20):}

\begin{table}[t]
\centering
\begin{tabular}{llll}
\toprule
Indicator & Score & Threshold & Theory \\
\midrule
global\_broadcast & 1.000 & 0.50 & GWT \\
hierarchical\_prediction & 1.000 & 0.40 & FEP \\
prediction\_error\_minimization & $0.889 \pm 0.038$ & 0.40 & FEP \\
sparse\_smooth\_coding & $0.834 \pm 0.057$ & 0.30 & FEP \\
global\_ignition\_nuanced & $0.725 \pm 0.031$ & 0.40 & GWT \\
agency & $0.698 \pm 0.009$ & 0.50 & FEP \\
ignition\_dynamics & $0.660 \pm 0.049$ & 0.40 & GWT \\
recurrent\_processing & 0.600 & 0.40 & GWT \\
attention\_control & 0.500 & 0.40 & AST \\
higher\_order\_representations & 0.500 & 0.50 & HOT \\
embodiment & 0.500 & 0.30 & AST \\
metacognition & 0.400 & 0.40 & HOT \\
bayesian\_binding\_quality & 0.300 & 0.30 & BLT \\
epistemic\_depth & 0.300 & 0.30 & BLT \\
\bottomrule
\end{tabular}
\caption{Passing indicators at warm start (14/20).}
\label{tab:passing-indicators}
\end{table}

\textbf{Failing (6/20):}

\begin{table}[t]
\centering
\small
\begin{tabular}{lllp{5.5cm}}
\toprule
Indicator & Score / Thresh. & Theory & Why It Fails \\
\midrule
attention\_schema & 0.250 / 0.50 & AST & Schema prediction requires longer tracking history \\
irreducibility & 0.300 / 0.40 & IIT & Decomposition limited without full neural connectivity \\
genuine\_implementation & 0.300 / 0.40 & BLT & Sustained recursive loops not yet established \\
integrated\_information & 0.200 / 0.30 & IIT & $\Phi$ approximation limited (PyPhi incompatibility) \\
local\_recurrence & 0.100 / 0.30 & RPT & Requires active substrates with feedback dynamics \\
algorithmic\_recurrence & 0.100 / 0.30 & RPT & Requires active substrates with feedback dynamics \\
\bottomrule
\end{tabular}
\caption{Failing indicators at warm start (6/20).}
\label{tab:failing-indicators}
\end{table}

The failing indicators share a common pattern: they measure features that require either extended operational history (attention schema prediction, genuine implementation), neural substrate connectivity data (integrated information, irreducibility, algorithmic recurrence, local recurrence), or both. These are the indicators most likely to improve when the framework is embedded in a system with persistent state and active neural substrates --- confirming the design intent described in Section~\ref{sec:embedding}.

\subsection{Ablation Study Results}
\label{sec:ablation}

The ablation study disables one theory module at a time and measures the impact on the overall score. Results from a single warm-start run:

\begin{table}[t]
\centering
\begin{tabular}{llll}
\toprule
Configuration & Overall Score & Drop from Baseline & Impact \\
\midrule
Baseline (all modules) & 0.513 & --- & --- \\
Without FEP & 0.368 & $-0.145$ (28.2\%) & Significant \\
Without GWT & 0.376 & $-0.136$ (26.6\%) & Significant \\
Without AST & 0.469 & $-0.043$ (8.5\%) & Moderate \\
Without HOT & 0.473 & $-0.040$ (7.7\%) & Moderate \\
Without BLT & 0.480 & $-0.032$ (6.3\%) & Moderate \\
\bottomrule
\end{tabular}
\caption{Ablation study results: impact of disabling individual theory modules.}
\label{tab:ablation}
\end{table}

Note: IIT and RPT are measurement-only modules --- they do not participate in the consciousness cycle as processing steps --- and therefore cannot be meaningfully ablated. Their indicators are assessed across the full system rather than by disabling a dedicated module.

\subsubsection{Interpretation}

\textbf{GWT and FEP are the most impactful modules.} Disabling either causes a 27--28\% drop in overall score. GWT provides the workspace competition and broadcast mechanism that all downstream modules depend on. FEP provides the prediction error signals and homeostatic drives that ground the system in embodied needs. Removing either disrupts the data flow to multiple other modules.

We note that this impact pattern is partly a consequence of the sequential pipeline architecture: GWT is the first processing phase, so all downstream modules depend on its output. A different module ordering might produce a different impact ranking. The ablation reveals architectural dependencies as much as theoretical importance, and we do not claim that GWT and FEP are more ``fundamental'' theories of consciousness --- only that they are more structurally central in our implementation.

\textbf{AST, HOT, and BLT have moderate but distinct impacts.} Removing any one causes a 6--8\% drop. These modules enrich the consciousness cycle --- adding self-modeling, meta-representation, and recursive inference --- but the cycle can still function without them.

\textbf{Cross-theory dependencies are visible.} When GWT is disabled, indicators belonging to HOT drop because workspace winners are no longer available for meta-representation. When FEP is disabled, BLT scores drop because the hierarchical predictive processor that BLT's hyper-model operates on is no longer producing prediction errors. These cross-module effects demonstrate that the theories interact in the multi-theory architecture rather than operating in isolation.

\subsubsection{Interaction Effects}

The ablation data reveals two notable interaction patterns:

\textbf{FEP $\rightarrow$ GWT coupling.} Prediction errors from the FEP module modulate candidate salience in the workspace competition (via cross-cycle feedback). When FEP is disabled, this salience boost disappears, and fewer candidates cross the ignition threshold. This supports the hypothesis that prediction error signals shape conscious access --- a cross-theory interaction that is invisible when either theory is implemented alone.

\textbf{GWT $\rightarrow$ HOT dependency.} When GWT is disabled, higher\_order\_representations drops to near zero. This is a direct architectural dependency: the HOT module generates meta-representations only for workspace winners. Without a functioning workspace, there are no winners to reflect on. This interaction is predicted by both theories --- GWT predicts that unconscious content does not receive global access, and HOT predicts that content without higher-order representation is not conscious.

Figure~\ref{fig:interactions} maps the cross-theory dependencies revealed by ablation.

% Figure 3: Cross-Theory Interaction Map
% Requires: tikz, arrows.meta, positioning, calc, decorations.pathmorphing
\begin{figure}[t]
\centering
\begin{tikzpicture}[
  % --- Global styles ---
  >={Stealth[length=5pt, width=4pt]},
  font=\sffamily,
  scale=0.92, every node/.style={scale=0.92},
  % --- Node styles by category ---
  structural/.style={
    circle, draw=#1, fill=#1!10, line width=1.6pt,
    font=\sffamily\bfseries, align=center, inner sep=0pt
  },
  structural/.default=blue!70!black,
  enrichment/.style={
    circle, draw=#1, fill=#1!8, line width=1.1pt,
    font=\sffamily\bfseries, align=center, inner sep=0pt
  },
  enrichment/.default=violet!60!black,
  measurement/.style={
    circle, draw=yellow!45!black, fill=yellow!4,
    line width=1pt, dashed, dash pattern=on 3pt off 2pt,
    font=\sffamily\bfseries, align=center, inner sep=0pt
  },
  integrative/.style={
    circle, draw=orange!65!black, fill=orange!10, line width=1.2pt,
    font=\sffamily\bfseries, align=center, inner sep=0pt
  },
  % --- Edge label style ---
  elabel/.style={
    font=\sffamily\scriptsize, fill=white, inner sep=1.5pt,
    outer sep=0.5pt, rounded corners=1pt
  },
]

% ============================================================
% NODE LAYOUT — hexagonal ring with GWT at top
%
% Positions chosen so edges route cleanly between nodes.
% Node sizes reflect ablation impact (larger = bigger drop).
% ============================================================

% --- GWT: top center ---
\node[structural, minimum size=30mm]
  (GWT) at (0, 3.0) {%
    {\large GWT}\\[1pt]
    {\fontsize{6.5}{7}\selectfont Workspace}\\[2pt]
    {\fontsize{6.5}{7}\selectfont\color{black!55} $-$26.6\%}
  };

% --- AST: upper left ---
\node[enrichment, minimum size=20mm]
  (AST) at (-3.3, 1.8) {%
    {\normalsize AST}\\[0pt]
    {\fontsize{5.5}{6}\selectfont Attention}\\[1pt]
    {\fontsize{5.5}{6}\selectfont\color{black!55} $-$8.5\%}
  };

% --- HOT: upper right ---
\node[enrichment, minimum size=19mm]
  (HOT) at (3.3, 1.8) {%
    {\normalsize HOT}\\[0pt]
    {\fontsize{5.5}{6}\selectfont Meta}\\[1pt]
    {\fontsize{5.5}{6}\selectfont\color{black!55} $-$7.7\%}
  };

% --- FEP: lower left (largest node) ---
\node[structural=red!60!black, minimum size=32mm]
  (FEP) at (-2.6, -0.6) {%
    {\large FEP}\\[1pt]
    {\fontsize{6.5}{7}\selectfont Inference}\\[2pt]
    {\fontsize{6.5}{7}\selectfont\color{black!55} $-$28.2\%}
  };

% --- BLT: lower right ---
\node[enrichment=orange!65!black, minimum size=18mm]
  (BLT) at (2.6, -0.6) {%
    {\normalsize BLT}\\[0pt]
    {\fontsize{5.5}{6}\selectfont Loop}\\[1pt]
    {\fontsize{5.5}{6}\selectfont\color{black!55} $-$6.3\%}
  };

% --- Damasio: bottom center ---
\node[integrative, minimum size=22mm]
  (Damasio) at (0, -2.6) {%
    {\fontsize{8.5}{9}\selectfont Damasio}\\[1pt]
    {\fontsize{5.5}{6}\selectfont Body}
  };

% --- IIT: far right, offset from main ring ---
\node[measurement, minimum size=14mm]
  (IIT) at (4.8, 3.4) {%
    {\fontsize{7.5}{8}\selectfont IIT}
  };

% --- RPT: far left, offset from main ring ---
\node[measurement, minimum size=14mm]
  (RPT) at (-4.8, 3.4) {%
    {\fontsize{7.5}{8}\selectfont RPT}
  };

% ============================================================
% DIRECTED EDGES — the 5 cross-theory interactions
% ============================================================

% --- 1. FEP -> GWT: prediction error -> salience (KEY FINDING) ---
%     Arc swings left of AST; label placed at explicit coordinates in clear space
\draw[->, line width=2.2pt, red!60!black]
  (FEP.120) .. controls (-5.0, 1.6) and (-3.4, 4.2) .. (GWT.195);
\node[elabel, text=red!60!black] at (-2.0, 3.85) {%
  \begin{tabular}{@{}c@{}}pred.\ error\\[-2pt]$\to$ salience\end{tabular}
};

% --- 2. GWT -> HOT: workspace winners -> meta-representation ---
\draw[->, line width=1.3pt, blue!50!black]
  (GWT.350) .. controls (1.8, 3.4) and (2.8, 2.8) .. (HOT.90)
  node[elabel, pos=0.45, above right=-1pt] {winners};

% --- 3. AST -> FEP: attention strength -> precision weights ---
\draw[->, line width=1.3pt, violet!55!black]
  (AST.250) -- (FEP.110)
  node[elabel, pos=0.45, right=3pt] {precision};

% --- 4. HOT -> BLT: higher-order thoughts -> epistemic depth ---
\draw[->, line width=1.3pt, violet!55!black]
  (HOT.300) -- (BLT.60)
  node[elabel, pos=0.50, right=2pt] {depth};

% --- 5. FEP -> Damasio: homeostatic drives -> protoself ---
\draw[->, line width=1.2pt, black!45]
  (FEP.290) .. controls (-2.0, -2.4) and (-1.4, -3.0) .. (Damasio.190)
  node[elabel, pos=0.40, left=2pt] {drives};

% ============================================================
% LEGEND — two rows, clean columns
% ============================================================
\begin{scope}[shift={(0, -4.4)}]
  \draw[gray!35, rounded corners=3pt, fill=gray!2]
    (-4.5, -0.15) rectangle (4.5, 0.95);

  % --- Row 1 (y=0.65) ---
  % Node size
  \fill[blue!10] (-4.0, 0.65) circle (4pt);
  \draw[blue!70!black, line width=1.2pt] (-4.0, 0.65) circle (4pt);
  \fill[blue!10] (-3.5, 0.65) circle (2.5pt);
  \draw[blue!70!black, line width=0.8pt] (-3.5, 0.65) circle (2.5pt);
  \node[font=\sffamily\fontsize{5.5}{6}\selectfont, anchor=west]
    at (-3.2, 0.65) {Size $=$ ablation impact};

  % Structural
  \fill[blue!10] (-0.5, 0.65) circle (3pt);
  \draw[blue!70!black, line width=1.2pt] (-0.5, 0.65) circle (3pt);
  \node[font=\sffamily\fontsize{5.5}{6}\selectfont, anchor=west]
    at (-0.2, 0.65) {Structural};

  % Enrichment
  \fill[violet!8] (1.2, 0.65) circle (3pt);
  \draw[violet!60!black, line width=0.9pt] (1.2, 0.65) circle (3pt);
  \node[font=\sffamily\fontsize{5.5}{6}\selectfont, anchor=west]
    at (1.45, 0.65) {Enrichment};

  % Integrative
  \fill[orange!10] (2.9, 0.65) circle (3pt);
  \draw[orange!65!black, line width=0.9pt] (2.9, 0.65) circle (3pt);
  \node[font=\sffamily\fontsize{5.5}{6}\selectfont, anchor=west]
    at (3.15, 0.65) {Integrative};

  % --- Row 2 (y=0.20) ---
  % Arrow width
  \draw[->, line width=2pt, red!60!black] (-4.1, 0.2) -- (-3.5, 0.2);
  \node[font=\sffamily\fontsize{5.5}{6}\selectfont, anchor=west]
    at (-3.2, 0.2) {Width $=$ interaction strength};

  % Measurement only
  \draw[yellow!45!black, dashed, dash pattern=on 2.5pt off 1.5pt,
        line width=0.8pt] (1.2, 0.2) circle (3pt);
  \node[font=\sffamily\fontsize{5.5}{6}\selectfont, anchor=west]
    at (1.5, 0.2) {Measurement only};

\end{scope}

\end{tikzpicture}

\caption{Cross-theory interaction map. Node size reflects ablation impact on the
overall consciousness score; percentages indicate the score decrease when that
module is disabled. Arrow thickness encodes interaction strength. The five
directed edges represent measurable cross-theory dependencies: FEP prediction
errors modulate GWT workspace salience (thickest arrow; the framework's key
finding), GWT workspace winners feed HOT meta-representations, AST attention
strength sets FEP precision weights, HOT outputs drive BLT epistemic depth,
and FEP homeostatic drives supply the Damasio protoself. IIT and RPT (dashed)
are measured across the system but do not participate as cycle steps and cannot
be ablated.}
\label{fig:interactions}
\end{figure}


\subsection{Reproducibility}
\label{sec:reproducibility}

All assessment results reported in this section can be reproduced by running \texttt{scripts/run\_assessment.py} on the released codebase:

\begin{lstlisting}[language=bash]
# Cold-start assessment
python scripts/run_assessment.py --report

# Warm-start assessment
python scripts/run_assessment.py --full --report

# Ablation study
python scripts/run_assessment.py --full --ablation
\end{lstlisting}

Cold-start results are perfectly deterministic. Warm-start results vary by $\pm 0.007$ (standard deviation across 10 runs, range 0.493--0.517) due to random synthetic stimuli during warmup. Passing indicator count showed zero variance: all 10 runs produced exactly 14/20 passing. All commands complete in under 5 seconds on consumer hardware.

The assessment does not require GPUs, network access, or external databases. It runs entirely on CPU using the framework's own modules. This is deliberate: reproducibility requires that no external dependencies affect the results.

\subsection{Noise Normalization and DCM}
\label{sec:noise-dcm-results}

The noise normalization methodology described in Section~\ref{sec:assessment} and the Digital Consciousness Model described in Section~\ref{sec:assessment} are implemented in the codebase and available for use. We do not report their results in this paper for two reasons: (1) noise normalization requires a separate characterization study to establish appropriate baseline noise models for each indicator, which is beyond the scope of this initial report; and (2) DCM scoring requires mapping each of the 13 perspectives to the framework's modules, which we have implemented but not yet validated against the perspectives' original criteria. Both are priorities for future work (Section~\ref{sec:future_work}). We include their methodology descriptions because they are part of the framework's assessment toolkit and will be used in subsequent publications.
